\documentclass[aspectratio=169]{beamer}

% --- Packages ---
\usepackage[utf8]{inputenc}
\usepackage{lmodern}
\usepackage{pifont}       % For checkmarks and crossmarks (\ding{51}, \ding{55})
\usepackage{tcolorbox}    % Modern styled callout boxes
\tcbuselibrary{skins,breakable}
\usepackage{booktabs}
\usepackage{tabularx}

% --- Beamer Theme & Palette ---
\usetheme{Madrid}
\usecolortheme{dove}
\setbeamertemplate{navigation symbols}{} % Remove default navigation icons

% --- Presentation Metadata ---
\title[EEL 5706 - HW2]{HW2 for EEL5706}
\author{Yousef Awad}
\date{\today}

\begin{document}

\begin{frame}{AN Code \& Residue Code Pros + Cons}

  \begin{columns}[T, onlytextwidth]

    % Column 1
    \begin{column}{0.48\textwidth}
      \centering{\textbf{\large AN Code}}
      \vspace{0.2cm}

      \begin{exampleblock}{Pros}
        \begin{itemize}\footnotesize
          \item Single Widened ALU path
          \item Immune to zero-residue masking via single-bit flips adding an error of $\pm 2$
        \end{itemize}
      \end{exampleblock}

      \begin{alertblock}{Cons}
        \begin{itemize}\footnotesize
          \item High decoding latency (via integer division)
          \item Complex multiplication overhead from $Ax * Ay = A^2xy$
        \end{itemize}
      \end{alertblock}
    \end{column}

    \hfill

    % Column 2
    \begin{column}{0.48\textwidth}
      \centering{\textbf{\large Residue Code}}
      \vspace{0.2cm}

      \begin{exampleblock}{Pros}
        \begin{itemize}\footnotesize
          \item Zero Decoding Latency
          \item Parallel check logic operates concurrently with execution path
        \end{itemize}
      \end{exampleblock}

      \begin{alertblock}{Cons}
        \begin{itemize}\footnotesize
          \item Requires hardware duplcation (residue ALU \& comparator logic)
          \item Vulnerable to zero-residue error masking during multiplication steps
        \end{itemize}
      \end{alertblock}
    \end{column}

  \end{columns}
\end{frame}

% --- Frame 2: Zero-Residue Problem & Mitigation ---
\begin{frame}{Mitigations to Zero-Residue Error Masking in Residue Code}

  % Top Section: Question & Context
  \begin{block}{\textbf{Original Question \& My Context of Understanding}}
    \small
    \textbf{Question:} \textit{For multiplication, Residue Code can encounter the
      zero-residue problem. That is, using $C_N*C_M$ to check $N*M$, if $C_N=0$, then any error that occurs in
      M cannot be detected. How to mitigate this problem? (25 pts, $C_N = N \bmod A$, $C_M = M \bmod A$)
    }
    \vspace{0.1cm}

    \textbf{My understanding of the problem:} \textit{Given that N is a multiple of A, and that $C_N = 0$ plus $C_{N*M} = C_N*C_M$,
      this means that $0 = (0*C_M) \bmod A$ which then expands to $0 = (0*[M \bmod A]) \bmod A$. $\therefore$ if
      any errors occur to M, and therefore $C_M$, it would go completely undetected.
    }
  \end{block}

  \vspace{0.15cm}

  % Bottom Section: Answer & Mitigations
  \begin{exampleblock}{\textbf{Mitigation Strategies}}
    \small
    \begin{itemize}\footnotesize
      \item \textbf{Operand Offset Shifting N:} Whenever $C_N = 0$ is detected, turn N to $N' = N + 1$ and then
        compute $(N + 1)*M - M = N*M$ so that M (and $C_M$) can correctly occur, and then on the final product it
        generates it subtracts M to get the correct result.
      \item \textbf{Inverse Residue Coding:} If you detect that $C_N = 0$ simply compute the compliment so
        that you can detect any faults in M or $C_M$. Therefore you would use $C'_N = A - (N \bmod A)$ to replace 
        $C_N$.
    \end{itemize}
  \end{exampleblock}

\end{frame}

\end{document}
